\uuid{Qqz0}
\titre{Matrice d'une application linéaire/Matrice de passage/Espace des matrices}
\auteur{hanani ; mzari}
\organisation{exo7}
\licence{CC BY-NC-SA 4.0}
\datecreate{2020-01-20}
\niveau{L1}
\module{Algèbre}
\chapitre{Matrice}
\sousChapitre{Matrice et application linéaire}
\difficulte{3}
\theme{Matrice d'une application linéaire, Matrice de passage, Espace des matrices}
\exo7qcmid{lille-2.matrice-application.026}

\contenu{
\begin{question}
\qclassification{108.03}
\qtags{motcle=Matrice d'une application linéaire/Matrice de passage/Espace des matrices}
On considère $M_2(\Rr)$ l'espace vectoriel des matrices carrées d'ordre $2$ à coefficients réels muni des deux bases ${\cal {B}}= \{ A_1,A_2,A_3,A_4\}$ et ${\cal {B'}}= \{ B_1,B_2,B_3,B_4\}$, où
$$A_1 = \left(\begin{array}{rc}1&0\\0&0\\ \end{array}\right) \; , \; A_2 = \left(\begin{array}{rc}0&1\\
0&0\\ \end{array}\right) \; , \; A_3 = \left(\begin{array}{rc}
0&0\\
1&0\\ \end{array}\right) \; , \; A_4 = \left(\begin{array}{rc}
0&0\\0&1\\ \end{array}\right) \; , $$
$$ B_1 = \left(\begin{array}{rc}1&1\\
0&0\\ \end{array}\right) \; , \; B_2 = \left(\begin{array}{rc}
0&1\\0&1\\ \end{array}\right) \; , \; B_3 = \left(\begin{array}{rc}
0&0\\ 1&1\\ \end{array}\right) \; ,\; B_4 = \left(\begin{array}{rc}
1&0\\-1&0\\ \end{array}\right).$$
On notera $P$ la matrice de passage de la base ${\cal {B}}$ à la base ${\cal {B'}}$ et $Q$ la matrice de passage de la base ${\cal {B'}}$ à la base ${\cal {B}}$.
\vskip1mm
Définition : Soit $E$ un espace vectoriel muni de deux bases ${\cal {B}}$ et ${\cal {B}}'$. La matrice de passage de la base ${\cal {B}}$ à la base  ${\cal {B}}'$ est la matrice de l'identité de $E$ de la base ${\cal {B}}'$ à la base  ${\cal {B}}$. Autrement dit, c'est la matrice dont la jième colonne est constituée des coordonnées du jième vecteur de la base ${\cal {B}'}$ dans la base  ${\cal {B}}$.
\vskip1mm
Quelles sont les assertions vraies ?
\begin{answers}
\good{$P = \left(\begin{array}{rccc}
1&0&0&1\\
1&1&0&0\\
0&0&1&-1\\
0&1&1&0\\
\end{array}\right).$}
\bad{$\displaystyle P = \frac{1}{2}\left(\begin{array}{rccc}
1&1&1&-1\\
-1&1&-1&1\\
1&-1&1&1\\
1&-1&-1&1\\
\end{array}\right).$}
\bad{$Q=\left(\begin{array}{rccc}
1&0&0&1\\
1&1&0&0\\
0&0&1&-1\\
0&1&1&0\\
\end{array}\right).$}
\good{La matrice de l'application identité de $M_2(\Rr)$ de la base ${\cal {B}}$ à la base ${\cal {B'}}$ est :
$\displaystyle \frac{1}{2}\left(\begin{array}{rccc}
1&1&1&-1\\
-1&1&-1&1\\
1&-1&1&1\\
1&-1&-1&1\\
\end{array}\right).$}
\end{answers}
\vskip2mm
\begin{explanations} $P = \left(\begin{array}{rccc}
1&0&0&1\\
1&1&0&0\\
0&0&1&-1\\
0&1&1&0\\
\end{array}\right)$ et
$\displaystyle Q= \frac{1}{2}\left(\begin{array}{rccc}
1&1&1&-1\\
-1&1&-1&1\\
1&-1&1&1\\
1&-1&-1&1\\
\end{array}\right)$. Par définition, $Q$ est la matrice de l'application identité de $M_2(\Rr)$ de la base ${\cal {B}}$ à la base ${\cal {B'}}$.
\end{explanations}
\end{question}
}
